title: Build System
alias: build
------------------------------

The original GDSiMS had a slightly cumbersome method to get a binary out of the source code - which is something we want to avoid. We want it to be as simple as possible to get a working binary on a new machine. As such, I've adopted a dual Makefile/CMake system that I have used well in the past, which allows a user to issue a single command, and have everything they need appear:

\begin{Code}
	cd /install/location/of/gdsims

	make
\end{Code}

This command alone should give users what they want: a compiled _gdsims_ (or _gdsims.exe_ on Windows) executable that they can then run. We also support a wide variety of other [[make options]].


\subsection{The Tiers of the Build System}

	There are three layers of the build system, designed to be used by different entities:

	* _Makefile_: a set of memorable shortcuts ([[make options]]), each of which names a preset. This is used by people compiling on the command line by calling _make_.
	* _CMakePresets.json_: the named build configurations (build directory, build type, and which options are on). Individuals using IDEs (such as Visual Studio) will use this\footnote{If its not detected automatically}, allowing an IDE to show the options in the configuration drop-down menu.
	* _CMakeLists.txt_: the actual build logic: fetching dependencies, running the precompiler-reflection\footnote{If installed, see [[external dependencies]]}, generating the [[version]] header, and compiling the library and the executable. Users should only use this if they want to perform manual cmake alterations; else it is read by the machine.
\subsection{Requirements}

	The build needs only:

	* A C++20 compiler (a recent GCC, Clang or MSVC)
	* CMake version 3.21 or later (needed for _CMakePresets.json_)
	* _git_ and an internet connection on the first build, so that CMake can fetch the [[external dependencies]]




